
\subsection{Bối cảnh}
Nông nghiệp hiện đại ngày càng sử dụng các cảm biến IoT để theo dõi các điều kiện môi trường như độ ẩm đất, nhiệt độ không khí, độ ẩm không khí, cường độ ánh sáng, độ pH và mực nước. Tuy nhiên, các nút cảm biến thường được bố trí trên diện tích lớn nên việc thu thập dữ liệu thủ công tốn nhiều thời gian và công sức.

Bên cạnh đó, một số khu vực nông nghiệp có kết nối Internet không ổn định. Việc lắp đặt gateway cố định hoặc hạ tầng truyền thông cho từng khu vực có thể làm tăng chi phí. Vì vậy, dữ liệu cảm biến có thể bị lưu riêng lẻ, trở nên lỗi thời hoặc khó quản lý tập trung.

\subsection{Giải pháp đề xuất}
Đề tài xây dựng nền tảng IoT có UAV hỗ trợ. UAV đóng vai trò như một gateway di động, tiếp cận các nút cảm biến để thu thập dữ liệu. Khi không có Internet, gateway có thể lưu dữ liệu cục bộ và thực hiện đồng bộ với hệ thống trung tâm khi kết nối được khôi phục.

Hệ thống hướng tới các chức năng chính:
\begin{itemize}
    \item Thu thập dữ liệu từ các nút cảm biến phân tán.
    \item Lưu trữ dữ liệu khi không có kết nối Internet.
    \item Đồng bộ dữ liệu với hệ thống trung tâm.
    \item Theo dõi dữ liệu hiện tại và dữ liệu lịch sử.
    \item Phát hiện giá trị bất thường và tạo cảnh báo.
\end{itemize}

